\begin{pastexam}[label={px:ccu-115-swa-03}]{中正資工所甲組 115 學年度 第 13 題}
（複選；全對才給分）

\begin{center}
\begin{tikzpicture}[x=1.45cm,y=1.45cm,>={Stealth[length=2.2mm]},
  v/.style={circle,draw,thick,minimum size=8mm,inner sep=0pt,font=\small\bfseries},
  e/.style={->,thick},
  l/.style={font=\small\bfseries,inner sep=1.5pt}]
% ---- Residual network G_f ----
\begin{scope}
\node[anchor=west,font=\large] at (-0.6,1.75) {Residual network G$_{\mathrm f}$};
\node[v] (s)  at (0,0)    {S};
\node[v] (v1) at (1.5,1)  {v1};
\node[v] (v2) at (1.5,-1) {v2};
\node[v] (v3) at (3,1)    {v3};
\node[v] (v4) at (3,-1)   {v4};
\node[v] (t)  at (4.5,0)  {t};
\draw[e] (v3) -- node[l,above]{12} (v1);
\draw[e,transform canvas={shift={(-0.045,0.068)}}] (s) -- node[l,above,sloped]{4} (v1);
\draw[e,transform canvas={shift={(0.045,-0.068)}}] (v1) -- node[l,below,sloped]{12} (s);
\draw[e,transform canvas={shift={(0.045,0.068)}}] (v2) -- node[l,above,sloped]{11} (s);
\draw[e,transform canvas={shift={(-0.045,-0.068)}}] (s) -- node[l,below,sloped,pos=0.4]{2} (v2);
\draw[e,transform canvas={xshift=-0.08cm}] (v1) -- node[l,left,rotate=90,anchor=south]{10} (v2);
\draw[e,transform canvas={xshift=0.08cm}] (v2) -- node[l,right,rotate=90,anchor=north]{4} (v1);
\draw[e] (v3) -- node[l,below right,pos=0.6]{9} (v2);
\draw[e] (v3) -- node[l,rotate=-90,anchor=south]{7} (v4);
\draw[e,transform canvas={shift={(0.045,0.068)}}] (v3) -- node[l,above,sloped]{1} (t);
\draw[e,transform canvas={shift={(-0.045,-0.068)}}] (t) -- node[l,below,sloped]{19} (v3);
\draw[e] (t) -- node[l,below,sloped]{4} (v4);
\draw[e,transform canvas={yshift=0.08cm}] (v4) -- node[l,above]{11} (v2);
\draw[e,transform canvas={yshift=-0.08cm}] (v2) -- (v4);
\end{scope}
% ---- Flow network G ----
\begin{scope}[yshift=-5.6cm]
\node[anchor=west,font=\large] at (-0.6,1.75) {Flow network G};
\node[v] (s)  at (0,0)    {S};
\node[v] (v1) at (1.5,1)  {v1};
\node[v] (v2) at (1.5,-1) {v2};
\node[v] (v3) at (3,1)    {v3};
\node[v] (v4) at (3,-1)   {v4};
\node[v] (t)  at (4.5,0)  {t};
\draw[e] (s) -- node[l,above,sloped]{12/16} (v1);
\draw[e] (s) -- node[l,below,sloped]{11/13} (v2);
\draw[e] (v1) -- node[l,above]{12/12} (v3);
\draw[e,transform canvas={xshift=-0.08cm}] (v1) -- node[l,left,rotate=90,anchor=south]{10} (v2);
\draw[e,transform canvas={xshift=0.08cm}] (v2) -- node[l,right,rotate=90,anchor=north]{4} (v1);
\draw[e] (v3) -- node[l,below right,pos=0.6]{9} (v2);
\draw[e] (v4) -- node[l,rotate=-90,anchor=south]{7/7} (v3);
\draw[e] (v3) -- node[l,above,sloped]{19/20} (t);
\draw[e] (v4) -- node[l,below,sloped]{4/4} (t);
\draw[e] (v2) -- node[l,below]{11/14} (v4);
\end{scope}
\end{tikzpicture}
\end{center}
(5\%) (\qquad) Consider the $S,t$ flow network shown above.

After three augmenting paths have already been applied, the residual graph contains a backward edge
with capacity 7.\\
Which of the following statements are \textbf{true}?\\
(a)~The current flow is maximum \qquad (b)~The flow violates capacity constraints\\
(c)~The residual graph must contain a cycle \qquad (d)~A previous augmenting path used that edge
forward\\
(e)~none of the above
\end{pastexam}

\begin{pastsoln}
答案：(a)(d)。

先核對兩張圖（原圖 $G_f$ 裡 v2$\to$v4 那條沒有標數字，由 $14 - 11 = 3$ 可知是 $3$）。 流量守恆：
\begin{center}
\small
\begin{tabular}{cll}
\toprule
頂點 & 流入 & 流出 \\
\midrule
v1 & S$\to$v1 $12$ & v1$\to$v3 $12$ \\
v2 & S$\to$v2 $11$ & v2$\to$v4 $11$ \\
v3 & v1$\to$v3 $12$ $+$ v4$\to$v3 $7$ $= 19$ & v3$\to$t $19$ \\
v4 & v2$\to$v4 $11$ & v4$\to$v3 $7$ $+$ v4$\to$t $4$ $= 11$ \\
\bottomrule
\end{tabular}
\end{center}
每條邊的流量都不超過容量，流量值是從 S 流出的 $12 + 11 = 23$。 剩餘容量照資結書定義~26.3：
原有的邊 $(u, v)$ 剩 $c(u, v) - f(u, v)$，反向剩 $f(u, v)$，算出來和 $G_f$ 一致。 「容量 $7$ 的
反向邊」是 v3$\to$v4，因為 v4$\to$v3 流了 $7$（而 v4$\to$v3 已經滿了，$G_f$ 裡沒有 v4$\to$v3）。

\begin{itemize}
\item (a) 對。 從 S 在 $G_f$ 裡走得到的頂點：S$\to$v1（$4$）、S$\to$v2（$2$）、v1$\leftrightarrow$v2、
v2$\to$v4（$3$）；v4 出去的只有 v4$\to$v2，v1 出去的 v1$\to$v3 已經滿了，v2$\to$v3 不存在（v3$\to$v2
沒有流量）。 所以走得到的是 $\{\mathrm{S}, \mathrm{v1}, \mathrm{v2}, \mathrm{v4}\}$，到不了 t，沒有
增廣路徑，目前的流是最大流（資結書定理~26.6）。 對應的割 $(\{\mathrm{S}, \mathrm{v1}, \mathrm{v2},
\mathrm{v4}\}, \{\mathrm{v3}, \mathrm{t}\})$ 的容量是 v1$\to$v3 $12$ $+$ v4$\to$v3 $7$ $+$ v4$\to$t $4$
$= 23$，等於流量值，再次確認。
\item (b) 錯。 上面逐邊檢查過，每條邊都 $0 \le f \le c$，而且守恆。
\item (c) 錯（照題意）。 「must」是說從題目的條件（有一條容量 $7$ 的反向邊）一定推得出有環，
這推不出來：反向邊 v3$\to$v4 本身不一定在環上，例如只有一條邊 $s \to t$ 而且流滿時，$G_f$ 只有
一條 $t \to s$，沒有環。 補充：就這張圖而言，$G_f$ 確實有環（例如 S$\to$v1$\to$S，因為 S$\to$v1
流了 $12$、還沒滿，兩個方向都有剩），如果把 (c) 讀成「這張 $G_f$ 有環」，它是對的。 我們照
「must」的語氣，把它當成一般性的推論，判為錯。
\item (d) 對。 v3$\to$v4 的剩餘容量 $7$ 等於 $f(\mathrm{v4}, \mathrm{v3}) = 7$。 原圖沒有 v3$\to$v4
這條邊，v4$\to$v3 的流量只能來自之前某條增廣路徑順著 v4$\to$v3 送流（資結書第~26.3~節）。
\item (e) 錯。
\end{itemize}
\end{pastsoln}
